\documentclass[a4paper, 12pt, oneside]{book}
\usepackage{amsmath, amssymb}
\usepackage{graphicx, enumerate, multirow, fancyhdr}
\usepackage{mathptmx, anyfontsize, setspace}
\usepackage{indentfirst}
\usepackage{algorithm}
\usepackage{algorithmic}
\usepackage{verbatim}
\usepackage{tikz}
\usetikzlibrary{positioning, shapes, arrows.meta}
\usepackage[total={5.5in, 9.5in},%
     top=1in, left=1.5in, includefoot]{geometry}%
\usepackage{tikz}
\usepackage{pgfplots}
\usepackage{float}
\usepackage{hyperref}
\usepackage{fancyhdr}
\usepackage{listings}
\usepackage{xcolor}
\usepackage{tocloft}
\usepackage{array}
\usepackage{tabularx}

% Graphics path for images
\graphicspath{{../Images/}}

\lstset{
  language=Python,
  basicstyle=\ttfamily\small,
  keywordstyle=\color{blue},
  commentstyle=\color{gray},
  stringstyle=\color{teal},
  numbers=left,
  numberstyle=\tiny\color{gray},
  stepnumber=1,
  numbersep=8pt,
  frame=single,
  breaklines=true,
  showstringspaces=false,
  tabsize=2
}


% Styles


% TikZ Libraries
\usetikzlibrary{shapes.geometric, arrows, positioning, fit, backgrounds, calc, shapes.multipart, decorations.pathreplacing,positioning,arrows.meta,shapes.multipart,decorations.pathreplacing}
\pgfplotsset{compat=1.17}

% Define colors
\definecolor{blockcolor}{RGB}{200,220,240}
\definecolor{arrowcolor}{RGB}{50,100,150}
\definecolor{componentcolor}{RGB}{180,200,220}
\definecolor{datacolor}{RGB}{255,230,200}

% Arrow styles
\tikzstyle{arrow} = [thick,->,>=stealth,arrowcolor]
\tikzstyle{block} = [rectangle, draw, fill=blockcolor, text width=6em, text centered, rounded corners, minimum height=3em]
\tikzstyle{component} = [rectangle, draw, fill=componentcolor, text width=8em, text centered, rounded corners, minimum height=4em]
\tikzstyle{data} = [cylinder, draw, fill=datacolor, shape border rotate=90, aspect=0.25, minimum height=3em, minimum width=4em]
%%
\setcounter{secnumdepth}{3}
%%
% Title of the project
\newcommand{\ttitle}[1]%
{\def\vtitle{#1}}%
% name of student
\newcommand{\nname}[1]%
{\def\vname{#1}}%
\newcommand{\rregisternumber}[1]%
{\def\vregisternumber{#1}}%

\newcommand{\nnamei}[1]%
{\def\vnamei{#1}}%
\newcommand{\rregisternumberi}[1]%
{\def\vregisternumberi{#1}}%

\newcommand{\nnameii}[1]%
{\def\vnameii{#1}}%
\newcommand{\rregisternumberii}[1]%
{\def\vregisternumberii{#1}}%

\newcommand{\nnameiii}[1]%
{\def\vnameiii{#1}}%
\newcommand{\rregisternumberiii}[1]%
{\def\vregisternumberiii{#1}}%
% Register Number of student
%\newcommand{\rregisternumber}[1]%
%{\def\vregisternumber{#1}}%
% Specialization
\newcommand{\sspecialization}[1]%
{\def\vspecialization{#1}}%
% File name of college emblem
\newcommand{\eemblem}[1]%
{\def\vemblem{#1}}%
% Department offering the MTech programme
\newcommand{\ddept}[1]%
{\def\vdept{#1}}%
% name of the College
\newcommand{\ccollege}[1]%
{\def\vcollege{#1}}%
% First line of address of College
\newcommand{\aaddresslinei}[1]%
{\def\vaddresslinei{#1}}%
% second line of address of College
\newcommand{\aaddresslineii}[1]%
{\def\vaddresslineii{#1}}%
% Third line of address of College
\newcommand{\aaddresslineiii}[1]%
{\def\vaddresslineiii{#1}}%
% Name of Head of Department
\newcommand{\hhod}[1]%
{\def\vhod{#1}}%
% Name of Internal Supervisor
\newcommand{\gguide}[1]%
{\def\vguide{#1}}%
%project co -ordinator1
\newcommand{\pprojectco}[1]%
{\def\pprojectco{#1}}%

%project co -ordinator2
\newcommand{\pprojectcoii}[1]%
{\def\pprojectcoii{#1}}%


% Date
\newcommand{\ddate}[1]%
{\def\vdate{#1}}%
%%
%% Redefining chapter/section/sub-section commands
%%
\makeatletter
\def\@makechapterhead#1%
{%
   \vspace*{30\p@}%
   {%
      \parindent \z@ \centering \normalfont
      \ifnum \c@secnumdepth >\m@ne
         \if@mainmatter
            \fontsize{15.5}{20}\selectfont\bfseries% 
            \@chapapp\space \thechapter
            \par\nobreak
         \fi
      \fi
      \interlinepenalty\@M
      \fontsize{15.5}{20}\selectfont\bfseries% 
      \MakeUppercase{#1}\par\nobreak
      \vskip 20\p@
   }%
}%
%%
\def\@chapter[#1]#2%
{%
   \thispagestyle{empty} 
   \ifnum \c@secnumdepth >\m@ne
      \if@mainmatter
         \refstepcounter{chapter}%
         \typeout{\@chapapp\space \thechapter.}%
         \addcontentsline{toc}{chapter}%
         {\protect\numberline{Chapter \thechapter.\enspace  
             \MakeUppercase{#1}}}%
      \else
         \addcontentsline{toc}{chapter}{#1}%
      \fi
   \chaptermark{#1}%
   \addtocontents{lof}{\protect\addvspace{10\p@}}%
   \addtocontents{lot}{\protect\addvspace{10\p@}}%
   \if@twocolumn
      \@topnewpage[\@makechapterhead{#2}]%
   \else
      \@makechapterhead{#2}%
      \@afterheading
   \fi
}%
%%
\def\@makeschapterhead#1%
{%
   \vspace*{30\p@}%
   {%
      \parindent \z@ \centering
      \normalfont
      \interlinepenalty\@M
      \fontsize{15.5}{20}\selectfont\bfseries  
      \MakeUppercase{#1}\par\nobreak
      \vskip 20\p@
  }%
}%
%%
\renewcommand\section{%
   \@startsection {section}{1}{\z@}%
   {-3.5ex \@plus -1ex \@minus -.2ex}%
   {2.3ex \@plus.2ex}%
   {\fontsize{13.5}{18}\selectfont\bfseries \MakeUppercase}%
}%
%%
\renewcommand\subsection{%
   \@startsection{subsection}{2}{\z@}%
   {-3.25ex\@plus -1ex \@minus -.2ex}%
   {1.5ex \@plus .2ex}%
   {\fontsize{12}{14}\selectfont\bfseries}%
}%
%%
\renewcommand\subsubsection{%
   \@startsection{subsubsection}{3}{\z@}%
   {-3.25ex\@plus -1ex \@minus -.2ex}%
   {1.5ex \@plus .2ex}%
   {\normalfont\normalsize}%
}%
%%
\renewcommand\thesubsubsection%
{%
\thesubsection.(\@roman\c@subsubsection)
}%
%%
\renewcommand{\l@section}{\@dottedtocline{1}{4.8em}{2em}}%
\renewcommand{\l@subsection}{\@dottedtocline{2}{9.6em}{2.8em}}%
%%
\makeatother
%%
\newcommand{\achapter}[1]%
{%
\chapter*{#1}%
\addcontentsline{toc}{chapter}{\MakeUppercase{#1}}%
}%
%%
\pagestyle{plain}
%%
%% KTU guidelines say that \parindent should be 
%% 10 character space. 10 character space is approximately 
%% equal to 3em (see the website indicated below).
%% A value of 2.8em has been chosen to align text  with
%% section headings.
%%https://www.microsoft.com/typography/developers/fdsspec/spaces.aspx#em_footnote
\setlength{\parindent}{2.8em}
%%